% FILE: 07_open_questions_and_next_work.tex
% Project: lifecycle
% Role: process-lifecycle-state-definitions-and-transitions

\section{Open Questions and Next Work}
\label{sec:lifecycle-open}

\subsection{Known Gaps}
\label{sec:lifecycle-open-gaps}

Four open questions are documented as known gaps in the lifecycle project.

\textbf{Gap 1: Executable test suite absent.} The checkpoint
\texttt{checkpoint\_\_lifecycle\_model\_complete.md} defines ALIVE as requiring
$\texttt{assert\_soundness}(N) \land \texttt{assert\_fitness}(N, L, 0.95) \land
\texttt{assert\_no\_ghost\_transitions}(N, L)$, but no executable test suite in
\texttt{lifecycle/} implements these assertions against real Petri nets or event logs.
The assertion blueprints are Python pseudocode specification artifacts. Until an
executable suite is authored and passes against \texttt{wasm4pm}-managed nets and
logs, the module-level ALIVE verdict cannot be issued. This is the primary gap between
the lifecycle project's PARTIAL status and an ALIVE status.

\textbf{Gap 2: wasm4pm execution obligations unresolved.} The MAPE-K mapping matrix
in \texttt{MAPE\_K\_MAP.md} lists eight capabilities required in \texttt{wasm4pm}
(online token replay, alignment computation, variant enumeration, repair plan evaluation,
actuation layer, receipt emitter, model store, metric time-series store) that are not
present in \texttt{wasm4pm-compat}. These obligations are extensively documented per
lifecycle phase, but \texttt{wasm4pm} is described as a separate downstream product.
The gap between lifecycle doctrine and \texttt{wasm4pm} execution authority is
explicitly unresolved. A crate that conflates the structure layer and the execution
layer is a structural defect, as stated in \texttt{MAPE\_K\_MAP.md}.

\textbf{Gap 3: Board Projection Holography formal connection absent.} The Board
Projection Holography (BPH) framework is introduced in
\texttt{board-projection-holography.md} as a paradigm for manufacturing board
projections directly from autonomic process ledgers. However, no formal connection
between BPH and the receipt chain data types or the \texttt{wasm4pm} types has been
established through a compiled, running bridge. The document references the
``Ostar Governor'' as the authority for mathematical closure of projections, but
the Ostar Governor is not formally defined within the lifecycle directory.

\textbf{Gap 4: PRL vs.\ Maturity Model canonicality unresolved.} The Process
Readiness Levels (PRL 0--5) defined in the readiness taxonomy parallel the five-level
Maturity Model (L1--L5) in \texttt{MATURITY\_MODEL.md} but differ in framing. PRL
characterizes the operational readiness of a specific process artifact; the Maturity
Model characterizes the governance maturity of a process as a whole. The relationship
between PRL classification and maturity level assignment, and which is canonical for
M\&A due diligence valuation purposes, is not explicitly resolved in any lifecycle
document.

\subsection{Deferred Gates}
\label{sec:lifecycle-open-deferred}

Three formal gates are deferred pending downstream work.

The \textbf{module-level ALIVE gate} is deferred pending (a) authorship of an
executable test suite implementing the four assertion blueprints, (b) a running
\texttt{wasm4pm} instance capable of managing Petri net artifacts and event logs,
and (c) passage of all four assertions against real nets and logs. The program-level
ALIVE checkpoint (\texttt{0x00}) does not substitute for this module-level gate.

The \textbf{MAPE-K loop closure gate} is deferred pending \texttt{wasm4pm} execution
engine implementation. The loop closure criterion requires that every Monitor
observation is a typed admitted artifact, every Analyze conclusion carries a bounded
confidence score, every Plan is typed and risk-scored, every Execute action produces a
receipt, and the Knowledge component can replay any past loop cycle. Condition 5
(historical replay from Knowledge) is the distinguishing criterion between a reactive
system and an autonomic one, and it requires the \texttt{wasm4pm} model store and
metric time-series store to be implemented.

The \textbf{Board Projection formal bridge gate} is deferred pending formal definition
of the projection engine's connection to the BLAKE3 receipt chain and the
\texttt{wasm4pm} typed artifact store. Until this bridge is specified and compiled,
BPH remains a research model (v30.1.1 designation) rather than an implemented
capability.

\subsection{Research Risks}
\label{sec:lifecycle-open-risks}

Three research risks are identified. [\textsc{Author Thesis / Interpretation}]

\textbf{Risk 1: Doctrine-execution divergence.} The lifecycle doctrine is specified
at high formality with Petri net mathematics, typed state machines, and formal gate
criteria. As \texttt{wasm4pm} is implemented, there is a risk that the execution
engine's type system will be unable to fully express the const generic constraints
(e.g., \texttt{WfNetConst<SOUNDNESS>}, \texttt{Between01<NUM, DEN>}) without nightly
Rust compiler features that may not stabilize on the required timeline. The intermediate
type alias workaround for E0391 variance cycles is already documented in the lifecycle
spec; additional workarounds may be required.

\textbf{Risk 2: Fitness threshold rigidity.} The 0.95 alignment fitness threshold at
Gate 3 is a declared board-level standard, not a derived empirical value. In real
enterprise processes, particularly those with high variance or seasonal patterns, a
0.95 threshold may trigger excessive $T_{\mathrm{elastic}}$ actuation without improving
governance. The threshold should be treated as a configurable parameter with formal
justification required for any value below 0.95. This is a risk to the lifecycle
framework's claim that Gate 3 is sufficient as stated.

\textbf{Risk 3: M\&A receipt chain scope.} The five JSON receipts in the parent
\texttt{receipts/} directory are explicitly noted as sample receipts requiring live
data for production use. If the lifecycle framework is deployed in an M\&A context
before production receipt infrastructure is in place, the sample receipts could be
mistaken for evidence of operational conformance. The RECEIPT\_REGISTRY's explicit
note that ``production receipts require live data'' must be preserved in all downstream
documentation.

\subsection{Next Receipted Work}
\label{sec:lifecycle-open-next}

The next receipted work items, in priority order, are as follows.

First, author an executable Python or Rust test suite implementing
\texttt{assert\_soundness}, \texttt{assert\_fitness}, \texttt{assert\_no\_ghost\_transitions},
and \texttt{assert\_decommission\_receipt} against at least one WF-net artifact and
one event log. Passage of all four assertions constitutes the lifecycle module-level
ALIVE gate and allows the PARTIAL status to be upgraded. This work does not require
\texttt{wasm4pm}; it requires only a pm4py or equivalent Python process mining library
and a well-formed WF-net.

Second, specify the \texttt{wasm4pm} capability surface required for online token
replay and receipt emission as a formal API contract (not prose description). This
specification should be expressed in the same typestate language as
\texttt{wasm4pm-compat} so that conformance between the specification and the
implementation is checkable at the type level.

Third, define the Board Projection Holography bridge as a formal type mapping from the
BLAKE3 receipt chain to the BPH projection engine inputs. The bridge should be
specifiable without a running \texttt{wasm4pm} instance, using the same
\texttt{wasm4pm-compat} types that the lifecycle framework already references.

Fourth, resolve the PRL vs.\ Maturity Model canonicality question by authoring a
single mapping document that assigns each PRL level to a Maturity Level and specifies
which is the canonical measure for M\&A due diligence valuation. This is a doctrine
document, not an implementation task, and can be completed without downstream
engineering work.
